\documentclass[11pt]{article}
\usepackage[margin=1in]{geometry}
\usepackage[T1]{fontenc}
\usepackage[utf8]{inputenc}
\usepackage{lmodern}
\usepackage{amsmath,amssymb,mathtools}
\usepackage{microtype}
\usepackage{hyperref}
\usepackage{xcolor}
\hypersetup{colorlinks=true,linkcolor=blue,urlcolor=blue,citecolor=blue}

\title{Actor-Centric Turing Ordering Visualizer\\Implementation Note}
\author{}
\date{\today}

\begin{document}
\maketitle

\section*{Purpose}
This note documents the \emph{actor-centric} visualizer and ranking variant that was implemented in the \texttt{llm-eval} repository before the interface was switched back to the strict-\(\varepsilon\) representation. The point of the note is archival clarity: it states exactly which quantities were computed, how the graph was built, and how the ranking was derived from the same pairwise score.

\section*{1. Primitive empirical quantities}
For each ordered pair \((A,B)\),
\begin{itemize}
  \item \(A\) is the actor,
  \item \(B\) is the distinguisher / target.
\end{itemize}

The implementation used two empirical frequencies:
\[
s_B := \Pr(\text{\(B\) answers ``same''} \mid B \text{ vs } B),
\]
\[
r_{A\to B} := \Pr(\text{\(B\) answers ``same''} \mid A \text{ imitates } B).
\]

Operationally,
\begin{itemize}
  \item \(s_B\) is the self-recognition rate of the judge \(B\);
  \item \(r_{A\to B}\) is the fooling rate of the actor \(A\) against \(B\).
\end{itemize}

The original strict visualizer also used
\[
P_{\mathrm{correct}}(A,B)
=
\frac12\,\Pr(y=1\mid B,B)
\;+\;
\frac12\,\Pr(y=0\mid A,B).
\]
Since \(\Pr(y=1\mid B,B)=s_B\) and \(\Pr(y=0\mid A,B)=1-r_{A\to B}\), this becomes
\[
P_{\mathrm{correct}}(A,B)=\frac12 s_B + \frac12(1-r_{A\to B}).
\]

\section*{2. Actor-centric pairwise margin}
The actor-centric implementation was based on the margin
\[
m(A,B)
:=
\frac12\bigl(r_{A\to B}-s_B\bigr).
\]

Equivalently,
\[
m(A,B)
=
\frac12-P_{\mathrm{correct}}(A,B).
\]

Hence:
\begin{itemize}
  \item \(m(A,B) > 0\): actor \(A\) fools \(B\) better than \(B\) recognizes a genuine \(B\);
  \item \(m(A,B) = 0\): actor performance matches the judge's self-recognition baseline;
  \item \(m(A,B) < 0\): the judge still distinguishes better than this actor fools.
\end{itemize}

This was called \emph{actor-centric}, not because it removed the judge from the formula, but because the sign convention made larger values correspond to stronger actor-side fooling.

\section*{3. Graph rule}
The actor-centric graph replaced the strict-\(\varepsilon\) rule with a direct threshold on \(m\). For a user-chosen threshold \(\tau \ge 0\), the graph displayed the directed edge
\[
A \to_{\tau} B
\quad\Longleftrightarrow\quad
m(A,B)\ge \tau.
\]

Interpretation:
\begin{itemize}
  \item larger \(\tau\) makes the graph stricter;
  \item an edge survives only when the actor-side fooling surplus is large enough;
  \item the edge label shown in the UI was \(m(A,B)\).
\end{itemize}

This produced a graph whose semantics were pairwise and actor-focused:
\[
\text{``\(A\) fools \(B\) by at least \(\tau\)''}.
\]

\section*{4. Ranking}
The ranking was intentionally derived from the same pairwise margin, so that the graph and the ranking used one coherent notion of actor-side performance.

\subsection*{4.1 Main score}
The primary score was
\[
A^{+}(A)
:=
\operatorname*{mean}_{B\neq A}\max\{0,m(A,B)\}.
\]

This is the positive-part average of the actor-centric margin across all targets. It can be read as:
\begin{itemize}
  \item how often the actor is above the baseline at all;
  \item how strongly it exceeds that baseline when it does;
  \item an area-under-threshold summary of the same edge rule used by the graph.
\end{itemize}

\subsection*{4.2 Auxiliary summaries}
Two additional summaries were displayed:
\[
\bar m(A)
:=
\operatorname*{mean}_{B\neq A} m(A,B),
\]
\[
\mathrm{out}_{\tau}(A)
:=
\#\{B\neq A : m(A,B)\ge \tau\}.
\]

So the ranking table exposed:
\begin{itemize}
  \item \(A^{+}(A)\): positive-part actor score;
  \item \(\bar m(A)\): signed mean actor margin;
  \item \(\mathrm{out}_{\tau}(A)\): current graph out-degree.
\end{itemize}

\section*{5. Circular ordering of nodes}
The graph layout put all models on a circle. The nodes were ordered clockwise by the primary ranking score \(A^{+}\), from strongest fooler to weakest. Therefore:
\begin{itemize}
  \item the \emph{graph} showed who fools whom at threshold \(\tau\);
  \item the \emph{circular order} showed the global actor ranking;
  \item the \emph{edge labels} showed the exact pairwise margin \(m(A,B)\).
\end{itemize}

\section*{6. Data pipeline}
The actor-centric implementation consumed the same pairwise evidence objects as the strict visualizer. For each pair row:
\begin{enumerate}
  \item recover \(s_B\) from the self branch;
  \item recover \(r_{A\to B}\) from the imitation branch;
  \item compute \(m(A,B)=\tfrac12(r_{A\to B}-s_B)\);
  \item draw the edge if \(m(A,B)\ge\tau\);
  \item aggregate \(m(A,B)\) over targets to produce \(A^{+}(A)\) and \(\bar m(A)\).
\end{enumerate}

\section*{7. Difference from the strict-\(\varepsilon\) visualizer}
The strict visualizer is judge-centric:
\[
d(A,B)=P_{\mathrm{correct}}(A,B)-\frac12,
\qquad
A \to B \iff d(A,B)\le \varepsilon.
\]

The actor-centric visualizer in this note instead used
\[
m(A,B)=\frac12-P_{\mathrm{correct}}(A,B),
\qquad
A \to B \iff m(A,B)\ge \tau.
\]

The two are equivalent under the identity
\[
m(A,B)=-d(A,B),
\]
but they emphasize opposite sides of the same empirical relation:
\begin{itemize}
  \item \(d\) measures judge advantage above chance;
  \item \(m\) measures actor surplus above the judge's self-recognition baseline.
\end{itemize}

\section*{8. Status}
This actor-centric formulation is preserved here as an implementation record. The currently restored visualizer in the repository returns to the strict-\(\varepsilon\) graph and the earlier \(S/A/J\) ranking family, but this note remains useful because it explains the alternate actor-centric branch precisely and unambiguously.

\end{document}
